\chapter{Cystoscopy}

\section*{What Is This Test or Procedure?}
Cystoscopy is an endoscopy of the urinary bladder via the urethra. It uses a rigid or flexible scope to inspect the bladder lining, urethra, and ureteral orifices.

\section*{Alternative Names}
Cystourethroscopy, Bladder scope

\section*{Principle}
A thin telescope with a lens and light source is passed through the urethra into the bladder, which is distended with sterile fluid. The bladder and urethral lining can be inspected directly, and instruments can be passed to biopsy lesions, remove stones, or treat abnormalities.
\section*{History}
Philipp Bozzini invented the first endoscope ('Lichtleiter') in 1806. Max Nitze built the first modern electric cystoscope in 1877.

\section*{Why This Test or Procedure Is Ordered}
Ordered to investigate hematuria (blood in urine), recurrent urinary tract infections, urinary incontinence, dysuria, or suspected bladder tumors.

\section*{Is This a Primary or Secondary Test or Procedure}
Primary diagnostic test for evaluating bladder pathology, urethral strictures, and hematuria.

\section*{How This Test or Procedure Is Performed}
The urethra is cleaned, and a local anesthetic gel is instilled. The cystoscope is inserted through the urethra into the bladder. Sterile water is used to fill and distend the bladder.

\section*{What Preparation Is Needed}
Empty your bladder before the test. You may be asked to take a prophylactic antibiotic.

\section*{Contraindications and Precautions}
Acute urinary tract infection (UTI) or acute prostatitis.

\section*{Risks}
Urinary tract infection, hematuria, temporary dysuria (burning during urination), or urinary retention due to urethral swelling.

\section*{What Happens After the Test or Procedure}
Drink plenty of water. You may feel a burning sensation during urination for 24-48 hours. Warm baths can help soothe discomfort.

\section*{Understanding Your Results}
Enables the urologist to detect stones, strictures, bladder cancer, or prostate enlargement.

\section*{Normal Results}
Smooth, normal urethra; normal prostate size; pink, normal bladder wall with no masses, stones, or inflammation.

\section*{Follow-up Tests and Procedures}
Bladder biopsy, urine cytology, or surgical intervention if abnormalities are found.

\section*{References}
\begin{enumerate}
  \item MedlinePlus. Cystoscopy. U.S. National Library of Medicine; 2025.
  \item Current Medical Diagnosis and Treatment. McGraw-Hill; 2025.
\end{enumerate}

\clearpage
